\clearpage
\begin{center}
{\large SGA 7 --- 第 XIV 讲}\par
\vspace{2em}
{\Large\underline{与超越理论的比较}}\par
\vspace{1.5em}
{\large P. Deligne 著}
\end{center}

\medskip
\noindent\textbf{原书目录}
\begin{center}
\begin{tabular}{@{}llr@{}}
引言 && 1\\
\S1. & 消灭圈的超越形式体系 & 1\\
1.1 & 穿孔圆盘的基本群：符号约定 & 2\\
1.2 & 集合层 & 5\\
1.3 & 函子 $R\Psi$ & 10\\
1.4 & 一个特殊情形 & 12\\
\S2. & 比较定理 & 16\\
\S3. & 孤立奇点 & 20\\
3.1 & 消灭圈簇 & 20\\
3.2 & Picard--Lefschetz 公式（超越情形） & 31\\
\S4. & de Rham 上同调 & 38
\end{tabular}
\end{center}

\medskip
\underline{引言}

为将第 XIII 讲中的消灭圈理论与其在 $\mathbb C$ 上的超越对应理论相比较，我们分两步进行。

\begin{enumerate}[label=(\alph*),leftmargin=2.6em]
\item 我们仿照第 XIII 讲建立一个超越形式体系，并用通常的方法证明比较定理（这些方法已用于证明 XIII, 2.3 中的有限性定理）。

\item 在孤立奇点的情形，我们更具体地说明这里引入的超越形式体系的含义。
\end{enumerate}

除比较定理外，本讲还包含超越情形下 Picard--Lefschetz 公式的一个证明。

在 \S4 中，我们给出消灭圈的 Hodge 理论的一项预备结果。本讨论班后续不会用到这一节。

\section*{1. 消灭圈的超越形式体系}

\subsection*{1.1. 穿孔圆盘的基本群：符号约定}

我们将考虑穿孔圆盘 $D^*$ 及其基本群 $\mathbb Z$。由于此群是交换群，可能出现符号歧义。我们通过忘掉它的交换性来消除这些歧义。

\subsubsection*{1.1.1.}
设 $X$ 是一个性质良好的连通拓扑空间，$x\in X$。
基本群的元素是以 $x$ 为基点的闭路的同伦类。
若 $\alpha,\beta\in\pi_1(X,x)$ 分别由闭路 $a,b$ 表示，则 $\alpha\beta$ 由先走 $b$ 再走 $a$ 所得的闭路表示。

\subsubsection*{1.1.2.}
设 $F$ 是 $X$ 上的局部常值层。对每条道路 $a:[0,1]\longrightarrow X$，$[0,1]$ 上的逆像 $a^*F$ 局部常值，因而是常值层，所以给出双射
\[
a:F_{a(0)}
\xleftarrow{\ \sim\ }H^0([0,1],a^*F)
\xrightarrow{\ \sim\ }F_{a(1)}.
\]
此双射只依赖于 $a$ 的同伦类。通过这一构造，$\pi_1(X,x)$ 在 $F_x$ 上左作用。函子 $F\longmapsto F_x$ 是从 $X$ 上局部常值层的范畴到带有 $\pi_1(X,x)$ 作用的集合的范畴的等价。

\subsubsection*{1.1.3.}
当 $x$ 变化时，群 $\pi_1(X,x)$ 是某个局部常值群层 $\Pi_1(X)$ 的纤维，此群层作用于每个局部常值层 $F$。若 $\pi_1(X,x)$ 是交换群，则 $\Pi_1(X)$ 是常值层；其值记为 $\pi_1(X)$。

\subsubsection*{1.1.4.}
群 $\pi_1(X,x)$ 作用于万有覆盖
\[
p:(\widetilde X,\widetilde x)\longrightarrow(X,x).
\]
将此覆盖视为 $X$ 上的平展空间，从而视为一个层（其局部截面层），则在 1.1.2 的对应下，它对应于带有特出元素 $e$ 的集合 $\pi_1(X,x)$，$\pi_1(X,x)$ 通过左平移作用于此集合。此群还通过 $\pi_1$-集合的自同构作用于这个 $\pi_1$-集合：元素 $\sigma\in\pi_1(X,x)$ 的作用是
\[
g\longmapsto g\sigma^{-1}.
\]
这一构造给出等同
\[
\pi_1(X,x)=\operatorname{Aut}_X(\widetilde X).
\]

若 $F$ 在 $X$ 上局部常值，则双射
\begin{equation}\tag{1.1.4.1}
H^0(\widetilde X,p^*F)
\longrightarrow(p^*F)_{\widetilde x}=F_x
\end{equation}
与 $\pi_1(X,x)$ 的作用交换：左端的作用是通过其在 $\widetilde X$ 上的作用搬运结构而得，右端则是 1.1.2 的作用。因此，1.1.2 的等价可以重新表述为：
\[
F\longmapsto
\bigl(\widetilde X\longmapsto
H^0(\widetilde X,p^*F)\bigr)
\]
是从 $X$ 上局部常值层的范畴到下述函子范畴的等价：
\[
(\text{万有覆盖，底空间为 }X)\longrightarrow(\mathrm{Sets}).
\]

\subsubsection*{1.1.5.}
若 $\pi_1(X,x)$ 是交换群，从而不依赖于 $x$，则 1.1.4 中 $\pi_1(X)$ 在 $\widetilde X$ 上的作用不依赖于基点：$\pi_1(X)$ 作用于 $X$ 的每个万有覆盖 $\widetilde X$。

必须记住，这个作用（也是我们将使用的唯一作用）对应于 1.1.3 中作用的逆。
